% ===== บทสรุปผู้บริหาร =====
\chapter*{บทสรุปผู้บริหาร}
\addcontentsline{toc}{chapter}{บทสรุปผู้บริหาร}

กลุ่มแผนงานใต้ร่มพระบารมี ดำเนินงานปีงบประมาณ 2569 ภายใต้ \textbf{3 โครงการหลักตามแบบ ง8}
กรอบงบประมาณ \textbf{8,000,000~บาท} มีรายการโครงการที่ดำเนินการจริงตลอดปี \textbf{75~รายการ}
หน่วยงานร่วมดำเนินการ 10~หน่วยงาน หัวหน้าโครงการ 40~ท่าน
ครอบคลุมพื้นที่ดำเนินงาน 48~พื้นที่ ในศูนย์พัฒนาโครงการหลวง สถานีเกษตรหลวง
โครงการตามพระราชดำริ และชุมชนบนพื้นที่สูง

\begin{tipbox}
\textbf{ผลสำเร็จที่โดดเด่น --- การเบิกจ่ายปิดปีสูงกว่าเป้าหมาย}\par
เบิกจ่ายทั้งปี \textbf{7,446,794~บาท คิดเป็นร้อยละ~93 ของกรอบงบประมาณ}
สูงกว่าเป้าหมายร้อยละ~90 ที่ตั้งไว้ในรายงานรอบ 10~เดือน
เมื่อนับรวมเงินที่ได้รับอนุมัติกันเหลื่อมปี 462,200~บาท จะคิดเป็น \textbf{ร้อยละ~98.9 ของกรอบ}
ไม่มีใบสั่งซื้อค้างชำระ และคืนงบประมาณเพียง 172,105~บาท
\end{tipbox}

\begin{infobox}
\textbf{ภาพการเร่งรัดในช่วงสองเดือนสุดท้าย}\par
ณ 4 สิงหาคม 2569 การเบิกจ่ายอยู่ที่ 4,460,332~บาท (ร้อยละ~56) และรายงานรอบ 10~เดือน
ระบุว่าต้องเบิกจ่ายเพิ่มอีกประมาณ 2.7~ล้านบาทภายใน 55~วัน
ผลจริงคือเบิกจ่ายเพิ่มได้ \textbf{2,986,462~บาท} ในช่วงเวลาดังกล่าว
ซึ่งมากกว่าที่ประเมินไว้ และทำให้ทุกโครงการหลักปิดปีเหนือร้อยละ~90
\end{infobox}

\begin{warnbox}
\textbf{ประเด็นที่ยังต้องแก้ในปีถัดไป}\par
(1) ยังไม่มีการบันทึกผลรายกิจกรรมเข้าระบบติดตาม (0 จาก 258~กิจกรรม)
ตัวเลขตัวชี้วัดในรายงานนี้จึงสะท้อน\textbf{ค่าที่รับมาดำเนินการ} มิใช่ผลที่ผ่านการตรวจรับ\par
(2) ตัวชี้วัดเชิงเครือข่ายและเชิงพาณิชย์ยังต่ำกว่าร้อยละ~20 ได้แก่ ตัวชี้วัดที่ 10 (18\%)
ที่ 17 (13\%) และที่ 38 (6\%)\par
(3) การเบิกจ่ายยังกระจุกตัวในช่วงปลายปีงบประมาณอย่างชัดเจน
\end{warnbox}

\section*{ภาพรวมตัวเลขสำคัญ}

\begin{table}[H]
\centering
\caption{สรุปผลการดำเนินงานภาพรวม ปิดปีงบประมาณ 2569 (ณ 30 กันยายน 2569)}
\small
\begin{tabular}{L{7.2cm}rr}
\toprule
\rowcolor{headerblue}
\textbf{รายการ} & \textbf{จำนวน} & \textbf{ร้อยละ}\\
\midrule
กรอบงบประมาณแผนงาน (บาท) & 8,000,000 & 100\%\\
\rowcolor{rowblue}
โอนงบประมาณให้ผู้รับผิดชอบแล้ว (บาท) & 7,992,353 & 99.9\%\\
\textbf{เบิกจ่ายจริง (บาท)} & \textbf{7,446,794} & \textbf{93.1\%}\\
\rowcolor{rowblue}
กันเงินเหลื่อมปีไปปีงบประมาณ 2570 (บาท) & 462,200 & 5.8\%\\
\textbf{เบิกจ่ายรวมเงินกันเหลื่อมปี (บาท)} & \textbf{7,908,994} & \textbf{98.9\%}\\
\rowcolor{rowblue}
คืนงบประมาณ (บาท) & 172,105 & 2.2\%\\
ใบสั่งซื้อค้างชำระ (บาท) & 0 & 0\%\\
\midrule
\rowcolor{rowblue}
รายการโครงการที่ดำเนินการ & 75 & ---\\
ระยะเวลาปีงบประมาณ & 365 จาก 365~วัน & 100\%\\
\bottomrule
\end{tabular}
\end{table}

\begin{table}[H]
\centering
\caption{ผลการเบิกจ่ายรายโครงการหลัก ปิดปีงบประมาณ 2569}
\small
\begin{tabular}{L{5.5cm}rrrr}
\toprule
\rowcolor{headerblue}
\textbf{โครงการหลัก} & \textbf{รายการ} & \textbf{กรอบงบ} & \textbf{เบิกจ่าย} & \textbf{ร้อยละ}\\
\midrule
ง8-1 ผลักดันเทคโนโลยี & 23 & 2,000,000 & 1,836,240 & 92\%\\
\rowcolor{rowblue}
ง8-2 ขับเคลื่อนองค์ความรู้ & 15 & 2,000,000 & 1,894,933 & 95\%\\
ง8-3 พัฒนากำลังคน & 37 & 4,000,000 & 3,715,621 & 93\%\\
\midrule
\rowcolor{headerblue}
\textbf{รวมทั้งแผนงาน} & \textbf{75} & \textbf{8,000,000} & \textbf{7,446,794} & \textbf{93\%}\\
\bottomrule
\end{tabular}
\end{table}

\section*{ข้อเสนอต่อผู้บริหารสำหรับปีงบประมาณ 2570}
\begin{enumerate}[leftmargin=1.6em,itemsep=2pt]
  \item \textbf{บังคับใช้วินัยการรายงานผลรายกิจกรรม} --- กำหนดให้การเบิกจ่ายงวดถัดไปผูกกับ
        การบันทึกผลและหลักฐานในระบบติดตาม เพื่อให้ตัวชี้วัดที่รายงานเป็นผลที่ตรวจรับได้จริง
  \item \textbf{กระจายการเบิกจ่ายรายไตรมาส} --- กำหนดเพดานการเบิกจ่ายขั้นต่ำรายไตรมาส
        เพื่อไม่ให้งานกระจุกตัวในสองเดือนสุดท้ายเหมือนปีนี้
  \item \textbf{ออกแบบตัวชี้วัดให้ตรงลักษณะงาน} --- ทบทวนตัวชี้วัดที่ 10 17 และ 38
        ซึ่งไม่สอดคล้องกับงานพัฒนาชุมชนบนพื้นที่สูง หรือออกแบบโครงการที่มุ่งตัวชี้วัดเหล่านี้โดยเฉพาะ
  \item \textbf{สนับสนุนกระบวนการพัสดุรายหน่วยงาน} --- หน่วยงานที่เบิกจ่ายต่ำ
        (มทร.ล้านนา ตาก ร้อยละ~52) มีข้อจำกัดด้านกระบวนการ มิใช่ด้านเนื้องาน
  \item \textbf{ผูกเป้าหมายการพัฒนาที่ยั่งยืนรายโครงการ} --- บันทึก SDG รายโครงการ
        ตั้งแต่ขั้นเสนอโครงการ เพื่อรายงานผลระดับสากลได้โดยไม่ต้องวิเคราะห์ย้อนหลัง
\end{enumerate}
